% ============================================================
% 05-conclusion.tex
% ============================================================
\section{Conclusion and Future Work}
\label{sec:conclusion}

This study investigated how centralized supervision and representation-specific evidence paths can support heterogeneous university counseling. Before production route repair, raw supervisor outputs matched domain specialists with 97.0\% agent accuracy while maintaining bounded execution in the evaluated cases. On the redesigned 100-query held-out benchmark, the full retrieval configuration achieved the highest Hit@1 ($0.7500$), Hit@3 ($0.9100$), Precision@5 ($0.2380$), Recall@5 ($0.8458$), and MRR@10 ($0.8354$) among the evaluated configurations. Its largest Hit@1 differences relative to E4 occurred on entity-comparison, temporal, and cross-domain queries, although E4 retained slightly higher cross-domain Recall@5. The E4-to-E5 comparison reflects the integrated routing, graph-grounding, and parent--child mapping bundle rather than an isolated component effect. In end-to-end evaluation, the full system achieved the highest Fact~EM and Source AP and tied for the highest Source Recall, but did not obtain the highest value on every LLM-judged metric. The ablation results therefore indicate metric-specific contributions rather than uniform superiority of the complete configuration.

Taken together, these findings suggest that no single representation or retrieval strategy dominates all university-counseling tasks. Relational curricula, structured financial information, and narrative regulations impose different evidence and reasoning requirements and are therefore better handled through separate but coordinated processing paths. In practice, the proposed architecture manages these paths within a unified deployment while restricting each specialist to its relevant knowledge sources and tools. The results support the integrated architecture in the evaluated setting, but do not establish that any individual component is universally optimal.

The evaluation remains limited to Can Tho University data, automated judging, offline knowledge schemas, small temporal and adversarial strata, and isolated robustness tests. The single-domain analysis also excludes the benchmark's 20 cross-domain queries. Future work will expand the held-out benchmark using anonymized real-world student counseling logs, subject to appropriate consent and privacy safeguards, to improve coverage of naturally occurring intents and conversational variation. Further work will evaluate transferability across institutions, conduct user studies with students and advisors, perform controlled computational-efficiency evaluation and model compression, employ cross-model evaluation to reduce judging bias, and compare \system{} with broader educational AI systems. It will also expand the production tool registry and evaluate matched single- and multi-agent configurations across increasing tool-set sizes, including production tools rather than synthetic distractors. Within the scope of this study, the central takeaway is not that one retrieval method should replace all others, but that institutional counseling may be better supported when evidence acquisition and execution constraints are aligned with the structure of each domain.
